\documentclass[11pt,a4paper]{article}
\usepackage[T1]{fontenc}
\usepackage[utf8]{inputenc}
\usepackage{lmodern}
\usepackage[margin=2.5cm]{geometry}
\usepackage{booktabs}
\usepackage{tabularx}
\usepackage{enumitem}
\usepackage{listings}
\usepackage{xcolor}
\usepackage[hidelinks]{hyperref}

\lstset{
  basicstyle=\ttfamily\small,
  columns=fullflexible,
  keepspaces=true,
  frame=single,
  rulecolor=\color{gray!50},
  backgroundcolor=\color{gray!8},
  breaklines=true,
  upquote=true,
  showstringspaces=false,
  xleftmargin=0.5em,
  xrightmargin=0.5em
}
\setlength{\emergencystretch}{3em}
\newcommand{\kbd}[1]{\texttt{#1}}
\newcolumntype{L}{>{\raggedright\arraybackslash}X}

\title{LeRobot (Python, English names) Installation Manual\\[0.3em]\large Installing LeRobot on another Mac}
\author{}
\date{10 November 2026}

\begin{document}
\maketitle

This manual explains how to install LeRobot, the Python drawing robot, on a
Mac where it has never been installed. This is the edition with English command names
(\kbd{fd()}, \kbd{qtr()}, \kbd{pd()}, \dots), in the folder
\kbd{LeRobot\_Python\_EN}, used by the book \emph{First Programming Lessons in
Python}; the French edition is in \kbd{LeRobot\_Python\_FR}. Once it is done, use the
\emph{LeRobot (Python, English names) User Manual} (\kbd{LeRobot\_Manual.pdf}) to draw and
to run robot programs.

The Python version needs much less than the OCaml version: no XQuartz, no
compiler, no startup task. The robot window is an ordinary macOS window
drawn with tkinter, the graphics library that comes with Python. The
installation takes about ten minutes. You need:
\begin{itemize}
\item a Mac with a recent macOS (Apple Silicon or Intel);
\item an administrator account (the Python installer asks for the Mac
  password);
\item an Internet connection;
\item a copy of the \kbd{LeRobot} folder, or at least its
  \kbd{LeRobot\_Python\_EN} folder (USB key, AirDrop, cloud drive, \dots).
\end{itemize}

Commands are typed in the \textbf{Terminal} application
(\emph{Applications} $\rightarrow$ \emph{Utilities} $\rightarrow$
\emph{Terminal}). Copy each command in one piece, paste it into the
Terminal, and press Return.

\tableofcontents

%======================================================================
\section{Python with Tk 8.6}

LeRobot needs Python 3.9 or later \emph{linked to Tk 8.6 or later}. Tk is the
library that draws the window.

\textbf{The \kbd{python3} that comes with macOS is not enough.} It is
linked to Tk 8.5, which no longer works on recent macOS: the robot window
does not open and the program hangs. Install Python from python.org, whose
installer includes a recent Tk:
\begin{enumerate}
\item Open \url{https://www.python.org/downloads/macos/} and download the
  \emph{macOS 64-bit universal2 installer} of the latest Python 3 release.
\item Open the downloaded \kbd{.pkg} file and follow the installer. At the
  end, the Finder window it opens can be closed.
\end{enumerate}
Check, in a new Terminal window:
\begin{lstlisting}[language=bash]
/Library/Frameworks/Python.framework/Versions/Current/bin/python3 -c "import tkinter; print(tkinter.TkVersion)"
\end{lstlisting}
The answer must be \kbd{8.6} or more.

You do not need to change which \kbd{python3} the Terminal uses: when
LeRobot is started with the macOS \kbd{python3}, it notices that its Tk is
too old and starts itself again with the python.org Python.

\paragraph{With Homebrew instead.} If Homebrew is installed
(\url{https://brew.sh}), its Python also works once its Tk module is
installed (replace 3.13 by the version that \kbd{brew info python} shows):
\begin{lstlisting}[language=bash]
brew install python@3.13 python-tk@3.13
\end{lstlisting}

%======================================================================
\section{The LeRobot folder}

Copy the \kbd{LeRobot} folder anywhere, for example onto the
\textbf{Desktop}. The Python programs are in its folder
\kbd{LeRobot\_Python\_EN}; the \kbd{LeRobotPythonEN.app} icon must stay in that
folder, next to \kbd{robot.sh}. The whole \kbd{LeRobot} folder can be moved
or renamed later: nothing in it depends on where it is.

\paragraph{Commands.} The \kbd{LeRobot} folder can be anywhere (on the
Desktop, in Documents, \dots): everything in it is found relative to it.
The commands of this manual are typed in a Terminal window \emph{opened in
the \kbd{LeRobot} folder}: type \kbd{cd} and a space, drag the
\kbd{LeRobot} folder from the Finder onto the Terminal window, and press
Return (or, in the Finder, right-click the folder and choose
\emph{Services} $\rightarrow$ \emph{New Terminal at Folder}).

Nothing has to be compiled. If the folder comes from another Mac, make the
scripts executable again (a copy can lose this):
\begin{lstlisting}[language=bash]
cd LeRobot_Python_EN && chmod +x robot.sh launch.sh check.sh
\end{lstlisting}

%======================================================================
\section{First launch}

\begin{enumerate}
\item Open the interactive drawing program:
\begin{lstlisting}[language=bash]
cd LeRobot_Python_EN
python3 Draw.py
\end{lstlisting}
  The robot window ($800 \times 700$ points) opens in front of the other
  windows. Press Return in it to finish, then \kbd{n} not to save.
\item Run a program of the book:
\begin{lstlisting}[language=bash]
./robot.sh book/03_programming/Square.py
\end{lstlisting}
  When the square is drawn, press Esc to close the window.
\item \kbd{LeRobotPythonEN.app}, in the \kbd{LeRobot\_Python\_EN} folder:
  double-click its icon to run Draw; drop a \kbd{.py} program on its
  icon to run it. The program runs in a Terminal window, which closes when
  the program ends normally and stays open if there is an error.

  The first time you open \kbd{LeRobotPythonEN.app}, macOS may refuse to open
  it because it comes from another Mac. Right-click the icon, choose
  \emph{Open}, then confirm. Allow it to control Terminal and to access the
  folder of \kbd{LeRobot} (Desktop, Documents, \dots) when macOS asks.
\item If the folder was copied from another Mac and the application still
  refuses to open, rebuild it on this Mac:
\begin{lstlisting}[language=bash]
cd LeRobot_Python_EN && tools/app/build.sh
\end{lstlisting}
\end{enumerate}

%======================================================================
\section{Troubleshooting}

\begin{tabularx}{\textwidth}{@{}LLL@{}}
\toprule
\textbf{Symptom} & \textbf{Cause} & \textbf{Fix} \\
\midrule
The program hangs and no window opens & Python linked to Tk 8.5 (the macOS \kbd{python3}) and no other Python found & Install Python from python.org (Section~1) \\
\addlinespace
\emph{il faut un Python avec Tk 8.6 ou plus} & Same cause: LeRobot found no suitable Python & Same fix \\
\addlinespace
\emph{No module named '\_tkinter'} & Homebrew Python without its Tk module & \kbd{brew install python-tk@3.13} (with the right version) \\
\addlinespace
\emph{No module named 'LibRobot'} & The program was started with \kbd{python3} from another folder & Use \kbd{./robot.sh Program.py}, or drop it on \kbd{LeRobotPythonEN.app} \\
\addlinespace
\emph{permission denied: ./robot.sh} & The scripts lost their executable mark during the copy & \kbd{chmod +x robot.sh launch.sh check.sh} (Section~2) \\
\addlinespace
\kbd{LeRobotPythonEN.app} does not open & macOS blocks applications copied from another Mac & Right-click, \emph{Open}; or rebuild it (Section~3) \\
\addlinespace
To use a given Python & --- & \kbd{LEROBOT\_PYTHON=\allowbreak /path/\allowbreak to/\allowbreak python3 \allowbreak ./robot.sh Program.py} \\
\bottomrule
\end{tabularx}

\end{document}
